%!TEX root = ../thesis.tex
% ******************************* Thesis Appendix B ********************************

\chapter{Quantum Metropolis Sampling}
\graphicspath{{Appendix2/Figs/}}


\section{Proving Approximate Quantum Detailed Balance of acceptance branch}
\label{appendix:Aproof}
In this section we prove that:
\begin{equation}
    \|\mathcal{A}^*_{\sigma_G}[\mathbb{I}] - \mathcal{A}[\sigma_G]\|_* = O\!\left(\frac{\beta\mu}{Nt}\right)\|\mathcal{A}[\sigma_G]\|_* + O\!\left(\epsilon_{\Delta}\right)
\end{equation}
\textbf{Truncate.} We begin with $\mathcal{A}$ and perform a truncation to $\tilde{\mathcal{A}}$ by substituting $\tilde \alpha, \tilde \beta$ for $\hat \alpha, \hat \beta$. Under this substitution $G^{pq}_j$ becomes $\tilde G^{pq}_j$ such that $\|\tilde G^{pq}_j\| =0$ for all $|p-q\pm\frac{N}{2\pi}E_jt|_{N} > \mu$.  By Lemma~\ref{lemma:0} this produces a modified operator $\widetilde{\text{BSI-QPE}}$ such that $\|\widetilde{\text{BSI-QPE}}- \text{BSI-QPE}\| \leq 2\epsilon_{\Delta}$. In Appendix~\ref{section:appendixTrunc} we prove the following Lemma that relates an operator norm bound approximation of a Kraus operator to a diamond-norm bound in the approximation of the CP-map.
\begin{lemma}
\label{lemma:3}
Let $\{D_s\}_{s\in S}$ be mutually orthogonal projectors satisfying $\sum_{s\in S} D_s \leq \mathbb{I}$. Given unitaries $U$ and $V$, define the measurement $K_s = U^{\dagger}D_s V$ and CP map $\mathcal{K}[\rho] = \sum_{s \in S} K_s\rho K_s^{\dagger}$. Further, given operators $\tilde U, \tilde V$, define approximate measurement operators $\tilde K_s = \tilde U^{\dagger}D_s \tilde V$ and the approximate map $\tilde{\mathcal{K}}[\rho] = \sum_{s \in S} \tilde K_s\rho \tilde K_s^{\dagger}$. If $\|U-\tilde U\| \leq \epsilon<\frac14$ and $\|V-\tilde V\| \leq \epsilon<\frac14$, then 
\begin{equation}
    \|\mathcal{K} - \tilde{\mathcal{K}}\|_{\diamond} \leq 6\epsilon
\end{equation}
\end{lemma}
We can see that $\mathfrak{E}[\rho] =  E \rho E$, $\mathcal{P}_S[\rho] = \sum_{s\in S} P_s \rho P_s$, $\mathcal{Q}_S[\rho] = \sum_{s\in S} Q_s \rho Q_s$ and $\mathcal{L}[\rho] = \sum_{bd} L_{bd} \rho L^{\dagger}_{bd}$ are all CP maps of the form described in Lemma~\ref{lemma:3}. Here we have dropped the $ae$ superscript from $E$, $P_s$, $Q_s$ and $L_{bd}$ by postponing the corresponding classical conditioning until the end of the circuit: the outcomes are stored in additional pointer registers, and the later operations are controlled on those registers. This is only a proof technique and does not alter the persistent registers $R, F$ used by $E$, $P$, $L$ and $Q$. As $E$ and $P_s$ can be written as $\text{BSI-QPE}^{\dagger}\cdot D_s\cdot\text{BSI-QPE}$, $Q = C^{\dagger}\cdot\text{BSI-QPE}^{\dagger}\cdot W\cdot D_s\cdot W\cdot \text{BSI-QPE}\cdot C$ and $L = \text{BSI-QPE}^{\dagger}\cdot D_s\cdot W\cdot \text{BSI-QPE}\cdot C$ we can use the operator norm bound $\|\widetilde{\text{BSI-QPE}}- \text{BSI-QPE}\| \leq 2\epsilon_{\Delta}$ to see that:
\begin{equation}
    \|\mathfrak{E} - \tilde{\mathfrak{E}}\|_{\diamond} \leq 12\epsilon_{\Delta} \qquad\|\mathcal{P}_S - \tilde{\mathcal{P}}_S\|_{\diamond} \leq 12\epsilon_{\Delta} \qquad\|\mathcal{Q}_S - \tilde{\mathcal{Q}}_S\|_{\diamond} \leq 12\epsilon_{\Delta} \qquad\|\mathcal{L} - \tilde{\mathcal{L}}\|_{\diamond} \leq 12\epsilon_{\Delta}
    \label{eq:bound_on_component_maps}
\end{equation}
We can combine these bounds using the following Lemma which we also prove in Appendix~\ref{section:appendixTrunc}.
\begin{lemma}
\label{lemma:composition}
Let $\mathcal{M}_1, \ldots, \mathcal{M}_k$ be completely positive trace-non-increasing maps and let $\tilde{\mathcal{M}}_1, \ldots, \tilde{\mathcal{M}}_k$ be completely positive maps with $\|\mathcal{M}_i - \tilde{\mathcal{M}}_i\|_{\diamond} \leq \varepsilon$ for all $i$. Then
\begin{equation}
    \|\mathcal{M}_k \circ \cdots \circ \mathcal{M}_1 - \tilde{\mathcal{M}}_k \circ \cdots \circ \tilde{\mathcal{M}}_1\|_{\diamond} \leq (1+\varepsilon)^k - 1 \leq e k \varepsilon,
\end{equation}
the last inequality holding whenever $k\varepsilon \leq 1$. More generally, if instead $\|\mathcal{M}_i\|_{\diamond}, \|\tilde{\mathcal{M}}_i\|_{\diamond} \leq c$ for some $c \geq 1$ and $\|\mathcal{M}_i - \tilde{\mathcal{M}}_i\|_{\diamond} \leq \varepsilon_i$, without assuming the $\mathcal{M}_i$ are trace-non-increasing, then
\begin{equation}
    \|\mathcal{M}_k \circ \cdots \circ \mathcal{M}_1 - \tilde{\mathcal{M}}_k \circ \cdots \circ \tilde{\mathcal{M}}_1\|_{\diamond} \leq c^{k-1}\sum_{i=1}^{k}\varepsilon_i .
\end{equation}
\end{lemma}
We can write $\mathcal{A}[\rho] = \mathcal{L} \circ \mathcal{E}[\rho]$, and hence by Lemma~\ref{lemma:composition} and \eqref{eq:bound_on_component_maps} we have:
\begin{equation}
    \|\mathcal{A} - \tilde{\mathcal{A}}\|_{\diamond} \leq 24e \epsilon_{\Delta}
\end{equation}
By the generalised H\"older's inequality of Schatten norms ($\|ABC\|_s \leq \|A\|_p\|B\|_q\|C\|_r$ for $\frac{1}{s} = \frac{1}{p}+ \frac{1}{q} + \frac{1}{r}$) we have:
\begin{align}
    \|\mathcal{A}^*_{\sigma_G}[\mathbb{I}] - \tilde{\mathcal{A}}^*_{\sigma_G}[\mathbb{I}]\|_* &\leq \|\sigma_G^{\frac{1}{2}}\|_2\|\mathcal{A}^{\dagger}[\mathbb{I}] - \tilde{\mathcal{A}}^{\dagger}[\mathbb{I}]\|_{\infty} \|\sigma_G^{\frac{1}{2}}\|_2\\
    &\leq \operatorname{Tr}[\sigma_G]\|\mathcal{A}^{\dagger}[\mathbb{I}] - \tilde{\mathcal{A}}^{\dagger}[\mathbb{I}]\|
\end{align}
as $\mathcal{M}^*_{\sigma}[\mathbb{I}] = \sigma^{\frac{1}{2}}\mathcal{M}^{\dagger}[\mathbb{I}]\sigma^{\frac{1}{2}}$, $\operatorname{Tr}[\sigma] = \|\sigma^{\frac{1}{2}}\|_2^2$ and $\|\cdot\|_{\infty}= \|\cdot\|$. Further, by the duality between the trace and operator norms: 
\begin{align}
    \|\mathcal{A}^{\dagger}[\mathbb{I}]- \tilde{\mathcal{A}}^{\dagger}[\mathbb{I}]\| &= \sup_{\|X\|_* \leq 1} |\operatorname{Tr}[X^{\dagger}(\mathcal{A}^{\dagger}[\mathbb{I}]- \tilde{\mathcal{A}}^{\dagger}[\mathbb{I}]) ]|\\
    &= \sup_{\|X\|_* \leq 1}|\operatorname{Tr}[\mathcal{A}[X]- \tilde{\mathcal{A}}[X] ]|\\
    &\leq \|\mathcal{A}- \tilde{\mathcal{A}}\|_{\diamond}
\end{align}
Hence,
\begin{equation}
    \|\mathcal{A}^*_{\sigma_G}[\mathbb{I}] - \tilde{\mathcal{A}}^*_{\sigma_G}[\mathbb{I}]\|_* \leq \operatorname{Tr}[\sigma_G]\,\|\mathcal{A}- \tilde{\mathcal{A}}\|_{\diamond} \leq 24e\epsilon_{\Delta}
    \label{eq:AAtildeclose1}
\end{equation}
and directly:
\begin{equation}
    \|\mathcal{A}[\sigma_G] - \tilde{\mathcal{A}}[\sigma_G]\|_* \leq \|\mathcal{A}-\tilde{\mathcal{A}}\|_{\diamond} \leq 24e\epsilon_{\Delta}
    \label{eq:AAtildeclose2}
\end{equation}
As $\sigma_G = \frac{1}{Z}\sum_j e^{-\beta E_j} \ket{\psi_j}\bra{\psi_j}$ where $\ket{\psi_j}$ are the eigenstates of $H$ with eigenvalues $E_j$. We can expand $\tilde{\mathcal{A}}_{\sigma_G}$ as:
    \begin{equation}
    \tilde{\mathcal{A}}_{\sigma_G}[\rho] = \int Tr_{RFM}\!\left[\sum_{a,b,d,e}\textcolor{blue}{\tilde L^{ade}_b} \textcolor{red}{\tilde E_{ae}} \sigma_G^{\frac{1}{2}} (\rho\otimes \ket{0}\bra{0}_R \otimes \frac{\mathbb{I}_F}{2}\otimes \ket{0}\bra{0}_M) \sigma_G^{\frac{1}{2}}\textcolor{red}{\tilde E_{ae}}\textcolor{blue}{\tilde L^{ade\dagger}_b}\right] d\mu(C)
\end{equation}
\begin{multline}
    \bra{\psi_m}\tilde{\mathcal{A}}_{\sigma_G}[\ket{\psi_j}\bra{\psi_k}]\ket{\psi_l}=\frac{1}{2Z} \sum_{abde} e^{-\frac{\beta}{2}(E_j+E_k)}\textcolor{blue}{f_{abde}} \int \textcolor{blue}{\braket{\psi_m|C|\psi_j}\braket{\psi_k|C^{\dagger}|\psi_l}} d\mu(C)\\\sum_{p_1zq_1 q_2}Tr\left[\textcolor{blue}{\tilde G^{bz\dagger}_m\ket{d}\bra{d}\tilde G^{bp_1}_m} \textcolor{red}{\tilde G_j^{ap_1\dagger}\ket{e}\bra{e}\tilde G_j^{a0}\tilde G_k^{a0\dagger}\ket{e}\bra{e}\tilde G_k^{aq_1}} \textcolor{blue}{\tilde G^{bq_1\dagger}_l\ket{d}\bra{d}\tilde G^{bq_2}_l} \delta_{zq_2}\right]
    \label{eq:Arho}
\end{multline}
Similarly, using $Tr[X\mathcal{M}^*_{\sigma}[\rho]] = Tr[\mathcal{M}_{\sigma}[X]\rho]$, (i.e. exchanging the roles of $\ket{\psi_l}\bra{\psi_m}$ and $\ket{\psi_j}\bra{\psi_k}$ above)
\begin{multline}
   \bra{\psi_m}\tilde{\mathcal{A}}^*_{\sigma_G}[\ket{\psi_j}\bra{\psi_k}]\ket{\psi_l} =\frac{1}{2Z} \sum_{abde} e^{-\frac{\beta}{2}(E_l+E_m)}f_{abde} \int \braket{\psi_k|C|\psi_l}\braket{\psi_m|C^{\dagger}|\psi_j} d\mu(C)\\\sum_{p_1z q_1 q_2} Tr\left[\tilde G^{bz\dagger}_k\ket{d}\bra{d}\tilde G^{bp_1}_k\tilde G_l^{ap_1\dagger}\ket{e}\bra{e}\tilde G_l^{a0}\tilde G_m^{a0\dagger}\ket{e}\bra{e}\tilde G_m^{aq_1}\tilde G^{bq_1\dagger}_j\ket{d}\bra{d}\tilde G^{bq_2}_j  \delta_{zq_2}\right]
    \label{eq:A*rho}
\end{multline}
% We have:
% \begin{equation}
%     \tilde G^{bp\dagger}_{j}\ket{0}\bra{0} \tilde G^{bq}_{j} = \begin{pmatrix}
%         (-1)^{p+q}\tilde \alpha^*(p-b+\frac{N}{2\pi}E_{j}t)\tilde  \alpha(q-b+\frac{N}{2\pi}E_{j}t) &(-1)^{b+p}\tilde  \alpha^*(p-b+\frac{N}{2\pi}E_{j}t)\tilde  \beta(q-b+\frac{N}{2\pi}E_{j}t)  \\(-1)^{b+q}\tilde  \beta^*(p-b+\frac{N}{2\pi}E_{j}t)\tilde  \alpha(q-b+\frac{N}{2\pi}E_{j}t) & \tilde  \beta^*(p-b+\frac{N}{2\pi}E_{j}t)\tilde  \beta(q-b+\frac{N}{2\pi}E_{j}t)
%     \end{pmatrix}
% \end{equation} 
% \begin{equation}
%     \tilde G^{bp\dagger}_{j}\ket{1}\bra{1} \tilde G^{bq}_{j} = \begin{pmatrix}
%         \tilde \beta^*(p-b-\frac{N}{2\pi}E_{j}t)\tilde  \beta(q-b-\frac{N}{2\pi}E_{j}t) &  -(-1)^{b+q}\tilde  \beta^*(p-b-\frac{N}{2\pi}E_{j}t)\tilde \alpha(q-b-\frac{N}{2\pi}E_{j}t) \\
%         -(-1)^{b+p}\tilde  \alpha^*(p-b-\frac{N}{2\pi}E_{j}t)\tilde  \beta(q-b-\frac{N}{2\pi}E_{j}t) & (-1)^{p+q}\tilde  \alpha^*(p-b-\frac{N}{2\pi}E_{j}t)\tilde \alpha(q-b-\frac{N}{2\pi}E_{j}t)
%     \end{pmatrix}
% \end{equation} 

\textbf{Reverse Histories.} The adjoint map is given by running the round backwards. Therefore, we can attempt to relate a branch of the adjoint map to a different branch  of the original map by reversing the measurement outcomes. This succeeds up to the Boltzmann factor, which we invert in the \textbf{Reweight} step using \eqref{eq:fabtofba}. We begin by decomposing the adjoint map into branches:
\begin{equation}
    \tilde{\mathcal{A}}^*_{\sigma_G}[\rho] =  \sum_{abdez}\tilde{\mathcal{A}}^*_{abdez,\sigma_G}[\rho]
\end{equation}
where 
\begin{multline}
  \tilde{\mathcal{A}}^*_{abdez,\sigma_G}[\rho] = \frac{1}{2Z} \sum_{jklm}  e^{-\frac{\beta}{2}(E_l+E_m)}{f_{abde}} {\ket{\psi_m}}{\bra{\psi_j}}\rho{\ket{\psi_k}}{\bra{\psi_l}} \int {\braket{\psi_k|C|\psi_l}\braket{\psi_m|C^{\dagger}|\psi_j}} d\mu(C)\\\sum_{p_1q_1}Tr\left[ {\tilde G^{bz\dagger}_k\ket{d}\bra{d}\tilde G^{bp_1}_k} {\tilde G_l^{ap_1\dagger}\ket{e}\bra{e}\tilde G_l^{a0}\tilde G_m^{a0\dagger}\ket{e}\bra{e}\tilde G_m^{aq_1}} {\tilde G^{bq_1\dagger}_j\ket{d}\bra{d}\tilde G^{bz}_j } \right]
\end{multline}
By performing the following ordered relabelling:
\begin{equation}
    a \rightarrow b+z \qquad b\rightarrow a+z \qquad p_1 \rightarrow q_1+ z \qquad q_1 \rightarrow p_1 + z \qquad e \leftrightarrow d \qquad z \rightarrow N-z
\end{equation}
Due to the shift-invariance of $G^{pq}_j = G^{p+z,q+z}_j$ we get:
\begin{multline}
  \tilde{\mathcal{A}}^*_{b+N-z,a+N-z,e,d,N-z,\sigma_G}[\rho] = \frac{1}{2Z} \sum_{jklm}  e^{-\frac{\beta}{2}(E_l+E_m)}{f_{b+N-z,a+N-z,e,d}} {\ket{\psi_m}}{\bra{\psi_j}}\rho{\ket{\psi_k}}{\bra{\psi_l}} \\\int {\braket{\psi_k|C|\psi_l}\braket{\psi_m|C^{\dagger}|\psi_j}} d\mu(C)\sum_{p_1q_1} Tr\left[{\tilde G^{a0\dagger}_k\ket{e}\bra{e}\tilde G^{aq_1}_k} {\tilde G_l^{bq_1\dagger}\ket{d}\bra{d}\tilde G_l^{b,z}\tilde G_m^{b, z\dagger}\ket{d}\bra{d}\tilde G_m^{bp_1}} {\tilde G^{ap_1\dagger}_j\ket{e}\bra{e}\tilde G^{a0}_j } \right]
\end{multline}
These substitutions reverse the history of the measurement outcomes. So the initial energy estimate $(-1)^ea$ and the final energy estimate $(-1)^db$ are exchanged. We further require that $d\mu(C) = d\mu(C^{\dagger})$ so:
\begin{equation}
    \int  \bra{\psi_k}C\ket{\psi_{l}}\bra{\psi_m}C^{\dagger} \ket{\psi_j} d\mu(C) = \int  \bra{\psi_k}C^{\dagger}\ket{\psi_{l}}\bra{\psi_m}C \ket{\psi_j} d\mu(C)
\end{equation}
which gives:
\begin{multline}
 \bra{\psi_m}\tilde{\mathcal{A}}^*_{b+N-z,a+N-z,e,d,N-z,\sigma_G}[\ket{\psi_j}\bra{\psi_k}]\ket{\psi_l}  = \frac{1}{2Z} e^{-\frac{\beta}{2}(E_l+E_m)}{f_{b+N-z,a+N-z,e,d}} \\\int {\braket{\psi_k|C^{\dagger}|\psi_l}\braket{\psi_m|C|\psi_j}} d\mu(C)\sum_{p_1q_1} Tr\left[ {\tilde G_m^{b, z\dagger}\ket{d}\bra{d}\tilde G_m^{bp_1}} {\tilde G^{ap_1\dagger}_j\ket{e}\bra{e}\tilde G^{a0}_j } \tilde G^{a0\dagger}_k\ket{e}\bra{e}\tilde G^{aq_1}_k \tilde G_l^{bq_1\dagger}\ket{d}\bra{d}\tilde G_l^{b,z}\right]
 \label{eq:endofhistorya}
\end{multline}
The RHS is now a branch of the original map up to the Boltzmann factor, which we deal with once we have identified its support.\\
\textbf{Support.} Band-limiting means that only a narrow set of measurement outcomes contribute, and identifying this set is what allows us to reweight the Boltzmann factors. Due to the restricted support of $\tilde \alpha$ and $\tilde \beta$ we have that $\|\tilde G^{bp\dagger}_{j}\ket{d}\bra{d}\tilde G^{bq}_{j}\| = 0$ for $|p-q|_{N} > 2\mu$. Further, we have $\|\tilde G^{ab \dagger}_j \ket{e}\| = 0$ if $|a-b-(-1)^e\frac{N}{2\pi}E_{j}t|_{N} > \mu$. Hence, we can restrict the summations in Eq.~\eqref{eq:endofhistorya}  to the set $a,b,d,e,p_1,q_1,z \in \chi_{jklm}$ of summation indices satisfying
\begin{equation}
\begin{aligned}
    |a-(-1)^e\frac{N}{2\pi}E_jt|_{N} &\leq \mu  &|a-(-1)^e\frac{N}{2\pi}E_kt|_{N} &\leq \mu\\
       |b-z-(-1)^{d}\frac{N}{2\pi}E_mt|_{N}& \leq \mu&|b-z-(-1)^{d}\frac{N}{2\pi}E_lt|_{N} &\leq \mu \\
    |p_1|_{N}&\leq 2\mu&  |q_1|_{N}&\leq 2\mu \\
    |z-p_1|_{N}&\leq 2\mu&  |z-q_1|_{N}&\leq 2\mu \\
\end{aligned}
\label{eq:conds}
\end{equation}
Applying the modular triangle inequality separately along the two paths gives the following necessary conditions, which define a superset $\chi_{jklm}'\supset\chi_{jklm}$:
\begin{equation}
\label{eq:condJ}
\begin{aligned}
    \left|(-1)^db -(-1)^ea-\frac{Nt}{2\pi}(E_m-E_j)\right|_{N}&\leq 6\mu,\\
    \left|(-1)^db -(-1)^ea-\frac{Nt}{2\pi}(E_l-E_k)\right|_{N}&\leq 6\mu,
\end{aligned}
\end{equation}
As $|E_j| \leq \|H\|$, if we choose $t = \frac{1}{16\|H\|}$ and $\mu < \frac{N}{96\pi}$ so that $E_j$ satisfies $-\frac{1}{8t} + 3\frac{2\pi}{Nt}\mu < E_j < \frac{1}{8t} - 3\frac{2\pi}{Nt}\mu$, then we have that $-\frac{N}{4}+6\mu < \frac{Nt}{2\pi}(E_i -E_j) < \frac{N}{4}-6\mu$ for all $i$ and $j$. Hence, using \eqref{eq:conds} we can restrict the ranges of $(-1)^db-(-1)^ea$ found within $\chi'$ to $kN - \frac{N}{4}< (-1)^db-(-1)^ea < kN + \frac{N}{4}$.
We define $\kappa_{abde}$ to be the principal representative of $(-1)^db-(-1)^ea$ modulo $N$ lying in $[-\frac{N}{2},\frac{N}{2})$, 
\begin{equation}
    \kappa_{abde} = (-1)^db-(-1)^ea-kN \text{ for } kN - \frac{N}{2}< (-1)^db-(-1)^ea < kN + \frac{N}{2}
    \label{eq:kappadef}
\end{equation}
and hence we have $-\frac{N}{4}<\kappa_{abde} < \frac{N}{4}$ and $-\frac{N}{4}+6\mu < \frac{Nt}{2\pi}(E_i -E_j) < \frac{N}{4}-6\mu$ for all $i$ and $j$ so we can rewrite \eqref{eq:condJ} without the modular norm as:
\begin{equation}
    \begin{aligned}
    \left|\kappa_{abde}-\frac{Nt}{2\pi}(E_m-E_j)\right|&\leq 6\mu,\\
    \left|\kappa_{abde}-\frac{Nt}{2\pi}(E_l-E_k)\right|&\leq 6\mu,
\end{aligned}
\label{eq:kappaclose}
\end{equation}
Further, if we consider:
\begin{align}
    \left|\kappa_{b+N-z,a+N-z,ed} - \kappa_{baed}\right|_{N} &= |(N-z)((-1)^d-(-1)^e)|_{N}\\
    &\leq 2|z|_{N}
\end{align}
As $-\frac{N}{4}<\kappa_{baed} < \frac{N}{4}$ for all $b,a,e,d$ on support of $\chi'_{jklm}$ we can drop the modular norm and as $|z|_{N} \leq 4\mu$ we have:
\begin{equation}
    \left|\kappa_{b+N-z,a+N-z,ed} - \kappa_{baed}\right| \leq 8\mu
    \label{eq:kappap2identity}
\end{equation}
We choose $f_{abde}$ to be defined as:
\begin{equation}
    f_{abde} = \min(1,e^{-\beta\frac{2\pi}{Nt}\kappa_{abde}})
\end{equation}
Using \eqref{eq:kappaclose}, \eqref{eq:kappap2identity} and $\kappa_{baed}
= -\kappa_{abde}$ we can see that for $a,b,d,e,p_1,q_1,z \in \chi_{jklm}'$:
\begin{equation}
    \begin{aligned}
    e^{-\frac{\beta}{2}E_m}\sqrt{f_{b+N-z,a+N-z,e,d}} &= e^{-\frac{\beta}{2}E_m}\min(1,e^{-\frac{\beta}{2}(\frac{2\pi}{Nt}\kappa_{baed}+ 8\xi^{\prime m,abdez}_j)})\\
    &= e^{-\frac{\beta}{2}E_m}\min(1,e^{-\frac{\beta}{2}(E_j-E_m + 6\xi^{m,abde}_j+ 8\xi^{\prime m,abdez}_j)})\\
    &= \min\!\left(e^{-\frac{\beta}{2}E_m}, e^{-\frac{\beta}{2}(E_j + 6\xi^{m,abde}_j+ 8\xi^{\prime m,abdez}_j)}\right)\\
    &= e^{-\frac{\beta}{2}(E_j+ 6\xi^{m,abde}_j+ 8\xi^{\prime m,abdez}_j)} \min(e^{-\frac{\beta}{2}(E_m-E_j-6\xi^{m,abde}_j- 8\xi^{\prime m,abdez}_j)},1)\\
     &= e^{-\frac{\beta}{2}(E_j+ 6\xi^{m,abde}_j+ 8\xi^{\prime m,abdez}_j)} \min(e^{-\frac{\beta}{2}(\frac{2\pi}{Nt}\kappa_{abde}- 8\xi^{\prime m,abdez}_j)},1)\\
    &= e^{-\frac{\beta}{2}E_j}\sqrt{f_{abde}}(1+u^{abdez}_{jm})
\end{aligned}
\end{equation}
\aitext{where the last equality is by the continuity of $\min(1,e^{-\frac{\beta}{2}\frac{2\pi}{Nt}x})$ with $|\xi^{m,abde}_j|,|\xi^{\prime m,abde}_j| \leq \frac{2\pi}{Nt}\mu$
and $|u^{abde}_{jm}| = O(\frac{\beta \mu}{Nt})$.} We can similarly see:
\begin{equation}
\begin{aligned}
e^{-\frac{\beta}{2}(E_l+E_m)}f_{b+N-z,a+N-z,e,d}
&=(1+u_{jm}^{abdez})(1+u_{kl}^{abdez})e^{-\frac{\beta}{2}(E_j+E_k)}f_{abde}.
\end{aligned}
\label{eq:fabtofba}
\end{equation}
\textbf{Reweight.} As the RHS of \eqref{eq:endofhistorya} is only non-zero for $a,b,d,e,p_1,q_1,z \in \chi_{jklm} \subset \chi_{jklm}'$, we can apply \eqref{eq:fabtofba} to \eqref{eq:endofhistorya} to get:
\begin{multline}
\bra{\psi_m}\tilde{\mathcal{A}}^*_{b+N-z,a+N-z,e,d,N-z,\sigma_G}[\ket{\psi_j}\bra{\psi_k}]\ket{\psi_l} 
   = \frac{(1+u_{jm}^{abdez})(1+u_{kl}^{abdez})}{2Z}   e^{-\frac{\beta}{2}(E_j+E_k)}{f_{abde}}  \\\int {\braket{\psi_k|C^{\dagger}|\psi_l}\braket{\psi_m|C|\psi_j}} d\mu(C)\sum_{p_1q_1}Tr\left[ {\tilde G_m^{b, z\dagger}\ket{d}\bra{d}\tilde G_m^{bp_1}} {\tilde G^{ap_1\dagger}_j\ket{e}\bra{e}\tilde G^{a0}_j } \tilde G^{a0\dagger}_k\ket{e}\bra{e}\tilde G^{aq_1}_k \tilde G_l^{bq_1\dagger}\ket{d}\bra{d}\tilde G_l^{b,z}\right]
\end{multline}
Therefore, we have the exact relation
\begin{equation}
  \bra{\psi_m}\tilde{\mathcal{A}}^*_{b+N-z,a+N-z,e,d,N-z,\sigma_G}[\ket{\psi_j}\bra{\psi_k}]\ket{\psi_l} 
   =(1+u_{jm}^{abdez})(1+u_{kl}^{abdez})\bra{\psi_m}\tilde{\mathcal{A}}_{abdez,\sigma_G}[\ket{\psi_j}\bra{\psi_k}]\ket{\psi_l}.
\end{equation}
As $|u_{jr}^{abde}| \leq \epsilon_{\mathcal{A}} = O(\frac{\beta \mu}{Nt})$ by Lemma~\ref{lemma:1}  we have:
\begin{equation}
    \|\tilde{\mathcal{A}}^*_{\sigma_G}[\mathbb{I}] - \tilde{\mathcal{A}}[\sigma_G]\|_* \leq (2\epsilon_{\mathcal{A}}+ \epsilon_{\mathcal{A}}^2)\|\tilde{\mathcal{A}}[\sigma_G]\|_*
    \label{eq:Aoutput}
\end{equation}
Therefore, combining \eqref{eq:AAtildeclose1}, \eqref{eq:AAtildeclose2} and \eqref{eq:Aoutput} via the triangle inequality, we find $\mathcal{A}$ satisfies the hypothesis of Lemma~\ref{theorem:quantumapproximatedeatildedbalance} as:
\begin{equation}
    \|\mathcal{A}^*_{\sigma_G}[\mathbb{I}] - \mathcal{A}[\sigma_G]\|_* = O\!\left(\frac{\beta\mu}{Nt}\right)\|\mathcal{A}[\sigma_G]\|_* + O\!\left(\epsilon_{\Delta}\right)
\end{equation}
\section{Proving Approximate Quantum Detailed Balance of rejection branch}
\label{appendix:Bproof}
In this section we will prove that:
\begin{equation}
    \|\mathcal{B}^*_{\sigma_G}(n)[\mathbb{I}] - \mathcal{B}(n)[\sigma_G]\|_* = O\left(\frac{(n+2)\beta\mu}{Nt} \right)\|\mathcal{B}(n)[\sigma_G]\|_* +O\left(\max\left((n+2)^2\mu\epsilon_{\Delta},(n+2)^{\frac{5}{2}}\sqrt{\frac{\beta }{Nt}}\mu^{\frac{3}{2}}\right)\right) 
\end{equation}
 First, we will redefine $\mathcal{B}(n)$ so our later history reversal argument will work. For $n\geq 1$, let $\mathcal{B}^{(s_0)}(n)$ denote
the map obtained by replacing the first $Q^{ae}_0$ by
$Q^{ae}_{s_0}$. We have $E_{ae}=P^{ae}_1E_{ae}$ and
\begin{equation}
    P^{ae}_0Q^{ae}_1P^{ae}_1
    =P^{ae}_0(\mathbb I-Q^{ae}_0)P^{ae}_1
    =-P^{ae}_0Q^{ae}_0P^{ae}_1.
\end{equation}
Hence the two Kraus operators differ only by a sign, so $\mathcal{B}^{(0)}(n)=\mathcal{B}^{(1)}(n)=\mathcal{B}(n)$. So we can redefine the map
\begin{equation}
    \mathcal{B}(n)=\begin{cases}\mathcal{B}^{(0)}(0)  &  \text{ for } n =0
    \\\frac{1}{2}\sum_{s_0=0}^1 \mathcal{B}^{(s_0)}(n) &  \text{ for }n\geq 1 \end{cases}
\end{equation}
\textbf{Truncate.} For either $s_0\in\{0,1\}$, we have
compositions
\begin{equation}
    \mathcal{B}^{(s_0)}(n) = \mathcal{P}_{\{1\}}\circ\mathcal{Q}_{\{0,1\}}\circ\mathcal{P}_{\{0\}}\circ\cdots\circ\mathcal{Q}_{\{0,1\}}\circ\mathcal{P}_{\{0\}}\circ\mathcal{Q}_{\{s_0\}}\circ\mathfrak{E}
\end{equation}
Hence, by Lemma~\ref{lemma:3}, Lemma~\ref{lemma:composition} and averaging over $s_0$ we have:
\begin{equation}
    \|\mathcal B(n)-\tilde{\mathcal B}(n)\|_\diamond\leq 24e(2n+3)\epsilon_{\Delta}.
\end{equation}
when $12(2n+3)\epsilon_{\Delta} \leq 1$. For this to hold for all $n \leq n^*$ we must have $\mu = \Omega(\frac{1}{\nu}\log n^*)$. Therefore, following the same argument as for $\mathcal{A}$:
\begin{equation}
    \|\mathcal{B}(n)[\sigma_G] - \tilde{\mathcal{B}}(n)[\sigma_G]\|_* \leq 24e(2n+3)\epsilon_{\Delta},
    \qquad \|\mathcal{B}^*_{\sigma_G}(n)[\mathbb{I}] - \tilde{\mathcal{B}}^*_{\sigma_G}(n)[\mathbb{I}]\|_* \leq 24e(2n+3)\epsilon_{\Delta}.
    \label{eq:Bsrelation}
\end{equation}
We can expand $\tilde{\mathcal{B}}_{\sigma_G}(n)$ as:
% \begin{multline}
%     \tilde{\mathcal{B}}_{\sigma_G}(n)[\rho] = \frac12\int Tr_{RFM}\bigg[ \sum_{a,e,s_0,\ldots,s_n}\tilde P^{ae}_1\tilde Q^{ae}_{s_n}\tilde P^{ae}_0\ldots \tilde Q^{ae}_{s_1}\tilde P^{ae}_0\tilde Q^{ae}_{s_0}\tilde E_{ae} \\\sigma_G^{\frac{1}{2}}(\rho\otimes \ket{0}\bra{0}_R\otimes \frac{\mathbb{I}_F}{2}\otimes \ket{0}\bra{0}_M)\sigma_G^{\frac{1}{2}}  \tilde E_{ae}\tilde Q^{ae}_{s_0}\tilde P^{ae}_0\tilde Q^{ae}_{s_1}\ldots\tilde P^{ae}_0\tilde Q^{ae}_{s_n}\tilde P^{ae}_1\bigg]d\mu(C).
% \end{multline}
% Therefore:
\begin{multline}
    \bra{\psi_m}\tilde{\mathcal{B}}_{\sigma_G}(n)[\ket{\psi_j}\bra{\psi_k}]\ket{\psi_l} = \frac{e^{-\frac{\beta}{2}(E_j+E_k)}}{2Z}\int Tr\bigg[ \sum_{a,e,s_0,\ldots,s_n}(\ket{\psi_l}\bra{\psi_m}\otimes \mathbb{I}_{RM})\tilde P^{ae}_1\tilde Q^{ae}_{s_n}\tilde P^{ae}_0\ldots\tilde Q^{ae}_{s_1}\tilde P^{ae}_0\tilde Q^{ae}_{s_0}\tilde E_{ae} \\(\ket{\psi_j}\bra{\psi_k}\otimes \ket{0}\bra{0}_R\otimes \frac{\mathbb{I}_F}{2}\otimes \ket{0}\bra{0}_M)  \tilde E_{ae}\tilde Q^{ae}_{s_0}\tilde P^{ae}_0\tilde Q^{ae}_{s_1}\ldots\tilde P^{ae}_0\tilde Q^{ae}_{s_n}\tilde P^{ae}_1\bigg]d\mu(C).
\end{multline}
Equivalently, using $Tr[X\mathcal{M}^*_{\sigma}[\rho]] = Tr[\mathcal{M}_{\sigma}[X]\rho]$, (i.e. exchanging the roles of $\ket{\psi_l}\bra{\psi_m}$ and $\ket{\psi_j}\bra{\psi_k}$ above):
\begin{multline}
       \bra{\psi_m}\tilde{\mathcal{B}}^*_{\sigma_G}(n)[\ket{\psi_j}\bra{\psi_k}]\ket{\psi_l} = \frac{e^{-\frac{\beta}{2}(E_l+E_m)}}{2Z}\int Tr\bigg[ \sum_{a,e,s_0,\ldots,s_n}(\ket{\psi_j}\bra{\psi_k}\otimes \mathbb{I}_{RM})\textcolor{red}{\tilde P^{ae}_1}\textcolor{blue}{\tilde Q^{ae}_{s_n}}\textcolor{red}{\tilde P^{ae}_0}\ldots\textcolor{blue}{\tilde Q^{ae}_{s_1}}\textcolor{red}{\tilde P^{ae}_0}\textcolor{blue}{\tilde Q^{ae}_{s_0}}\textcolor{red}{\tilde E_{ae}} \\(\ket{\psi_l}\bra{\psi_m}\otimes \ket{0}\bra{0}_R\otimes \frac{\mathbb{I}_F}{2}\otimes \ket{0}\bra{0}_M)  \textcolor{red}{\tilde E_{ae}}\textcolor{blue}{\tilde Q^{ae}_{s_0}}\textcolor{red}{\tilde P^{ae}_0}\textcolor{blue}{\tilde Q^{ae}_{s_1}}\ldots\textcolor{red}{\tilde P^{ae}_0}\textcolor{blue}{\tilde Q^{ae}_{s_n}}\textcolor{red}{\tilde P^{ae}_1}\bigg]d\mu(C)
    \label{eq:Bstarexpanded}
\end{multline}
We can expand the operators in these expressions to get:
\begin{multline}
    \bra{\psi_m}\tilde{\mathcal{B}}_{\sigma_G}(n)[\ket{\psi_j}\bra{\psi_k}]\ket{\psi_l} = \frac{e^{-\frac{\beta}{2}(E_j+E_k)}}{2Z}\\\sum_{a,e,\{b\},\{c\},\{d\},\{d'\},\{f\},\{r\},z} \bar g^{ae}(\{b\},\{c\},\{d\},\{d'\}) D^{ae,z}_{jklm}(\{b\},\{c\},\{d\},\{d'\},\{f\},\{r\})I_{jklm}(\{f\},\{r\})
    \label{eq:def}
\end{multline}
\begin{multline}
    \bra{\psi_m}\tilde{\mathcal{B}}^*_{\sigma_G}(n)[\ket{\psi_j}\bra{\psi_k}]\ket{\psi_l} = \frac{e^{-\frac{\beta}{2}(E_l+E_m)}}{2Z}\\\sum_{a,e,\{b\},\{c\},\{d\},\{d'\},\{f\},\{r\},z} \bar g^{ae}(\{b\},\{c\},\{d\},\{d'\}) D^{ae,z}_{lmjk}(\{b\},\{c\},\{d\},\{d'\},\{f\},\{r\})I_{lmjk}(\{f\},\{r\})
    \label{eq:B*def}
\end{multline}
where the contribution from the $F$ qubit is given by:
\begin{multline}
    D^{ae,z}_{lmjk}(\{b\},\{c\},\{d\},\{d'\},\{f\},\{r\}) = 
    \sum_{\{p\},\{q\}}
Tr\bigg[\textcolor{red}{\tilde G^{az\dagger}_m\ket{e}\bra{e}\tilde G^{aq_1}_m}\textcolor{blue}{\tilde G^{c_1q_1\dagger}_{r_1}\ket{d'_1}\bra{d'_1}\tilde G^{c_1q_2}_{r_1}}\textcolor{red}{\tilde G^{c_2q_2\dagger}_{r_2}(\mathbb{I} - \delta_{ac_2}\ket{e}\bra{e})\tilde G^{c_2q_3}_{r_2}}\\\textcolor{blue}{\tilde G^{c_3q_3\dagger}_{r_3}\ket{d'_2}\bra{d'_2}\tilde G^{c_3q_4}_{r_3}\ldots
\tilde G^{c_{2n+1}q_{2n+1}\dagger}_{r_{2n+1}}\ket{d'_{n+1}}\bra{d'_{n+1}}\tilde G^{c_{2n+1}q_{2n+2}}_{r_{2n+1}}}\textcolor{red}{\tilde G^{aq_{2n+2}\dagger}_j\ket{e}\bra{e}\tilde G^{a0}_j}\\ \textcolor{red}{\tilde G^{a0\dagger}_k\ket{e}\bra{e}\tilde G^{ap_{2n+2}}_k}\textcolor{blue}{\tilde G^{b_{2n+1}p_{2n+2}\dagger}_{f_{2n+1}}\ket{d_{n+1}}\bra{d_{n+1}}\tilde G^{b_{2n+1}p_{2n+1}}_{f_{2n+1}}\ldots \tilde G^{b_3p_4\dagger}_{f_3}\ket{d_2}\bra{d_2}\tilde G^{b_3p_3}_{f_3}}\\\textcolor{red}{\tilde G^{b_2p_3\dagger}_{f_2}(\mathbb{I} - \delta_{ab_2}\ket{e}\bra{e})\tilde G^{b_2p_2}_{f_2}}\textcolor{blue}{\tilde G^{b_1p_2\dagger}_{f_1}\ket{d_1}\bra{d_1}\tilde G^{b_1p_1}_{f_1}}\textcolor{red}{\tilde G^{ap_1\dagger}_l\ket{e}\bra{e}\tilde G^{az}_l}\bigg]
\end{multline}
and the system register gives:
\begin{multline}
I_{lmjk}(\{f\},\{r\})=\int \textcolor{blue}{\bra{\psi_m}C^{\dagger}\ket{\psi_{r_1}}\bra{\psi_{r_1}}C\ket{\psi_{r_2}}\bra{\psi_{r_2}}C^{\dagger}\ket{\psi_{r_3}}\bra{\psi_{r_3}}C\ldots \bra{\psi_{r_{2n}}}C^{\dagger}\ket{\psi_{r_{2n+1}}}\bra{\psi_{r_{2n+1}}}C\ket{\psi_{j}}}\\\textcolor{blue}{
\bra{\psi_{k}}C^{\dagger}\ket{\psi_{f_{2n+1}}}\bra{\psi_{f_{2n+1}}}C\ket{\psi_{f_{2n}}}\ldots C^{\dagger}\ket{\psi_{f_3}}\bra{\psi_{f_3}}C\ket{\psi_{f_2}}\bra{\psi_{f_2}}C^{\dagger}\ket{\psi_{f_1}}\bra{\psi_{f_1}}C\ket{\psi_l}}d\mu(C)
\end{multline}
and the $M$ qubit gives:
\begin{equation}
       \bar g^{ae}(\{b\},\{c\},\{d\},\{d'\}) = \begin{cases} \textcolor{blue}{\bra{0}R_{0}^{ac_1d'_1e}R_{0}^{ab_1d_1e}\ket{0}}& \text{ for } n = 0\\
       \textcolor{blue}{\frac12\sum_{s_0,\ldots,s_n}
       \bra{0}R_{s_0}^{ab_{2n+1}d_{n+1}e}R_{s_1}^{ab_{2n-1}d_ne}\ldots
       R_{s_n}^{ab_{1}d_{1}e} R_{s_n}^{ac_{1}d'_{1}e}\ldots
       R_{s_1}^{ac_{2n-1}d'_ne}R_{s_0}^{ac_{2n+1}d'_{n+1}e}\ket{0}} & \text{ for } n \geq 1
       \end{cases}
       \label{eq:gbar-definition}
\end{equation}
\textbf{Reverse Histories.} We define $\tilde{\mathcal{B}}^*_{z,\sigma_G}(n)$ to be the term of \eqref{eq:B*def} with fixed final register value $z$, so that $\tilde{\mathcal{B}}^*_{\sigma_G}(n) = \sum_z \tilde{\mathcal{B}}^*_{z,\sigma_G}(n)$, and define $\tilde{\mathcal{B}}_{z,\sigma_G}(n)$ from \eqref{eq:def} in the same way. By applying the following ordered relabelling to $\tilde{\mathcal{B}}^*_{z,\sigma_G}(n)$:
\begin{align}
    a &\rightarrow a+ z &\qquad b_i &\rightarrow c_{2n+2-i} + z\qquad& c_i &\rightarrow b_{2n+2-i} + z  \qquad p_i \rightarrow q_{2n+3-i}+z \qquad q_i \rightarrow p_{2n+3-i} +z\\   d_i &\leftrightarrow d'_{n+2-i} &\qquad f_i &\leftrightarrow r_{2n+2-i} &\qquad z &\rightarrow N -z
\end{align}
We can see:
\begin{equation}
   D^{a+N-z,e,N-z}_{lmjk}(\{c+N-z\}^R,\{b+N-z\}^R,\{d'\}^R,\{d\}^R,\{r\}^R,\{f\}^R) =  D^{ae,z}_{jklm}(\{b\},\{c\},\{d\},\{d'\},\{f\},\{r\})
\end{equation}
\begin{equation}
    I_{lmjk}(\{r\}^R,\{f\}^R) = I_{jklm}(\{f\},\{r\})
\end{equation}
where $\{\cdot\}^R$ denotes reversal.
Relabelling the summed paths therefore gives
\begin{multline}
     \bra{\psi_m}\tilde{\mathcal{B}}^*_{N-z,\sigma_G}(n)[\ket{\psi_j}\bra{\psi_k}]\ket{\psi_l} = \frac{e^{-\frac{\beta}{2}(E_l+E_m)}}{2Z}\sum_{\substack{a,e,\{b\},\{c\},\\\{d\},\{d'\},\{f\},\{r\}}} \bar g^{a+N -z,e}(\{c+N -z\}^R,\{b+N -z\}^R,\{d'\}^R,\{d\}^R) 
     \\D^{ae,z}_{jklm}(\{b\},\{c\},\{d\},\{d'\},\{f\},\{r\})I_{jklm}(\{f\},\{r\}).
     \label{eq:relabelline}
\end{multline}
In Appendix~\ref{section:appendixTrunc} we prove the following lemma which will enable us to reverse the histories:
\begin{lemma}
\label{lemma:summed-history-reversal}
For $n\geq1$, let $\{x_i\}_{i=1}^{n+1}$ and $\{y_i\}_{i=1}^{n+1}$ denote
the compound labels
\begin{equation}
    y_i=(b_{2n+3-2i},d_{n+2-i}),\qquad x_i=(c_{2n+3-2i},d'_{n+2-i}),
\end{equation}
with $a,e$ fixed.  Choose $\theta_x\in[0,\pi/2]$ so that
\begin{equation}
    W_x=W_{bd}=\begin{pmatrix}
                \sqrt{1 - f_{abde}} & \sqrt{f_{abde}}\\
       \sqrt{f_{abde}} & -\sqrt{1 - f_{abde}}
    \end{pmatrix}= \begin{pmatrix}\cos\theta_x&\sin\theta_x\\
    \sin\theta_x&-\cos\theta_x\end{pmatrix},
\end{equation}
and set $\Delta_i=\theta_{y_i}-\theta_{x_i}$.  Then
\begin{equation}
    \bar g^{ae}(\{b\},\{c\},\{d\},\{d'\})
    =\frac12\cos\Delta_1\cos\Delta_{n+1}
      \prod_{i=1}^{n}\cos(\Delta_i-\Delta_{i+1}).
\end{equation}
Consequently,
\begin{equation}
    \bar g^{ae}(\{b\},\{c\},\{d\},\{d'\})
    =\bar g^{ae}(\{c\}^R,\{b\}^R,\{d'\}^R,\{d\}^R).
    \label{eq:PRR}
\end{equation}
\end{lemma}
Hence, by Lemma~\ref{lemma:summed-history-reversal} we can exchange $\{b\}^R \leftrightarrow \{c\}$ and $\{d\}^R \leftrightarrow \{d'\}$ within $g$ for $n\geq 1$, and as $R_0^{abde}$ is real and hermitian we can exchange $b \leftrightarrow c$ and $d \leftrightarrow d'$ when $n=0$:
\begin{multline}
         \bra{\psi_m}\tilde{\mathcal{B}}^*_{N-z,\sigma_G}(n)[\ket{\psi_j}\bra{\psi_k}]\ket{\psi_l} = \frac{e^{-\frac{\beta}{2}(E_l+E_m)}}{2Z}\\\sum_{\substack{a,e,\{b\},\{c\},\\\{d\},\{d'\},\{f\},\{r\}}} \bar g^{a+N -z,e}(\{b+N -z\},\{c+N -z\},\{d\},\{d'\}) 
         \\D^{ae,z}_{jklm}(\{b\},\{c\},\{d\},\{d'\},\{f\},\{r\})I_{jklm}(\{f\},\{r\}).  
    \label{eq:relabelline-first}
\end{multline}
Let us consider the CP-maps $\tilde{\mathcal{T}}^*_{N-z}(n)$ which is given by $\tilde{\mathcal{B}}^*_{N-z}(n)$ with every $W_{abde}$ replaced with $W_{a+z,b+z,de}$. Then we have:
\begin{multline}
         \bra{\psi_m}\tilde{\mathcal{T}}^*_{N -z,\sigma_G}(n)[\ket{\psi_j}\bra{\psi_k}]\ket{\psi_l} = \frac{e^{-\frac{\beta}{2}(E_l+E_m)}}{2Z}\\\sum_{a,e,\{b\},\{c\},\{d\},\{d'\},\{f\},\{r\}} \bar g^{a,e}(\{b\},\{c\},\{d\},\{d'\}) D^{ae,z}_{jklm}(\{b\},\{c\},\{d\},\{d'\},\{f\},\{r\})I_{jklm}(\{f\},\{r\}).
     \label{eq:relabelline-symmetrized}
\end{multline}

\textbf{Support.} We now consider the support of \eqref{eq:relabelline-symmetrized}. The restricted support of $\tilde\alpha$ and $\tilde\beta$ implies
$\|\tilde G^{bp\dagger}_{j}(\mathbb{I} - \delta_{ab}\ket{e}\bra{e})\tilde G^{bq}_{j}\| = 0$ and
$\|\tilde G^{bp\dagger}_{j}\ket{d}\bra{d}\tilde G^{bq}_{j}\| = 0$
for $|p-q|_{N} > 2\mu$, while
$\|\tilde G^{ap \dagger}_j \ket{e}\| = 0$ if
$|a-p-(-1)^e\frac{Nt}{2\pi}E_j|_{N}>\mu$.  Thus a nonzero
term must obey
\begin{equation}
\begin{aligned}
    |a-z-(-1)^e\tfrac{Nt}{2\pi}E_j|_{N} &\leq \mu
    &|a-z-(-1)^e\tfrac{Nt}{2\pi}E_k|_{N} &\leq \mu\\
    |c_{2i-1}-q_{2i-1}-(-1)^{d'_i}\tfrac{Nt}{2\pi}E_{r_{2i-1}}|_{N} &\leq \mu
    &|b_{2i-1}-p_{2i-1}-(-1)^{d_i}\tfrac{Nt}{2\pi}E_{f_{2i-1}}|_{N} &\leq \mu,
    \quad 1\leq i\leq n+1,\\
    |a-(-1)^{e}\tfrac{Nt}{2\pi}E_m|_{N} &\leq \mu
    &|a-(-1)^{e}\tfrac{Nt}{2\pi}E_l|_{N}& \leq \mu.\\
        |q_{2n+2}|_{N} &\leq 2\mu & |p_{2n+2}|_{N} &\leq 2\mu\\
    |q_{i+1}-q_i|_{N} &\leq 2\mu
    &|p_{i+1}-p_i|_{N} &\leq 2\mu,\quad 1\leq i\leq2n+1,\\
    |z-q_1|_{N} &\leq 2\mu & |z-p_1|_{N} &\leq 2\mu
\end{aligned}
\label{eq:conjE}
\end{equation}
Following either chain from $p_1$ or $q_1$ gives
\begin{equation}
    |p_i|_{N},|q_i|_{N} \leq (4n+4)\mu < (4n+6)\mu,
    \qquad 1\leq i\leq 2n+2. \qquad |z|_{N} \leq (4n+6)\mu
    \label{eq:pqbounds}
\end{equation}
Further, using the definition of $\kappa_{abde}$ \eqref{eq:kappadef} and \eqref{eq:conjE} we have:
\begin{align}
    \left|\kappa_{ab_{2i-1}d_i e}-\left(\tfrac{Nt}{2\pi}(E_{f_{2i-1}}-E_j)+(-1)^{d_i}p_{2i-1}\right)\right|_{N}&\leq 2\mu,\\
    \left|\kappa_{ac_{2i-1}d'_i e}-\left(\tfrac{Nt}{2\pi}(E_{r_{2i-1}}-E_k)+(-1)^{d'_i}q_{2i-1}\right)\right|_{N}&\leq 2\mu,
    \qquad 1\leq i\leq n+1.
    \label{eq:kappaebounds}
\end{align}
Since $|E_i|\leq\|H\|$ and $t=\frac{1}{16\|H\|}$, every scaled energy difference $\frac{Nt}{2\pi}(E_i - E_j)$ has absolute value at most $N/(16\pi)<\frac{N}{4}$. This taken together with \eqref{eq:pqbounds}-\eqref{eq:kappaebounds} gives:
\begin{equation}
    |\kappa_{ab_{2i-1}d_i e}|,|\kappa_{ac_{2i-1}d'_i e}| < \frac{N}{4}+(4n+6)\mu,
    \qquad 1\leq i\leq n+1.
    \label{eq:Qkappabound}
\end{equation}
These are the $\kappa$ values used in $\tilde{\mathcal{T}}^*_{N-z}(n)$. Applying the same argument directly to \eqref{eq:B*def} gives the same bound for the $\kappa_{abde}$ used in $\tilde{\mathcal{B}}^*_z(n)$. By considering the ends of the chains in \eqref{eq:conjE} and using \eqref{eq:pqbounds} we have:
\begin{equation}
\begin{aligned}
    \left|\frac{Nt}{2\pi}(E_m-E_j)\right|_{N}&\leq(4n+8)\mu,\\
    \left|\frac{Nt}{2\pi}(E_l-E_k)\right|_{N}&\leq(4n+8)\mu.
\end{aligned}
\end{equation}
Again since $|\frac{Nt}{2\pi}(E_i - E_j)| \leq \frac{N}{4}$ both modular norms may be replaced by ordinary absolute values:
\begin{equation}
\label{eq:energiessame}
\begin{aligned}
    \left|\frac{Nt}{2\pi}(E_m-E_j)\right|&\leq(4n+8)\mu,\\
    \left|\frac{Nt}{2\pi}(E_l-E_k)\right|&\leq(4n+8)\mu.
\end{aligned}
\end{equation}
Therefore,
\begin{equation}
\begin{aligned}
    \bra{\psi_m}\tilde{\mathcal{T}}^*_{N -z,\sigma_G}(n)[\ket{\psi_j}\bra{\psi_k}]\ket{\psi_l}=0\quad\text{unless}\quad
    \left|\frac{Nt}{2\pi}(E_m-E_j)\right|&\leq(4n+8)\mu,\\
    \left|\frac{Nt}{2\pi}(E_l-E_k)\right|&\leq(4n+8)\mu.
    \label{eq:Bsupport}
\end{aligned}
\end{equation}
\textbf{Reweight.} For $n\geq1$, define on this support (assuming $\mu \leq \frac{Nt}{(n^*+2)\beta}$):
\begin{equation}
    v_{jm}:=e^{-\frac{\beta}{2}(E_m-E_j)}-1=O(\frac{(n+2)\beta \mu}{Nt}).
\end{equation}
with $v_{jm}=0$ outside of this support.
Using $e^{-\frac{\beta}{2}(E_l+E_m)}=(1+v_{kl})(1+v_{jm})e^{-\frac{\beta}{2}(E_j+E_k)}$, we can rewrite \eqref{eq:relabelline-symmetrized} exactly as:
\begin{multline}
    \bra{\psi_m}\tilde{\mathcal{T}}^*_{N -z,\sigma_G}(n)[\ket{\psi_j}\bra{\psi_k}]\ket{\psi_l} =  \frac{(1+v_{kl})(1+v_{jm})e^{-\frac{\beta}{2}(E_j+E_k)}}{2Z}\\\sum_{a,e,\{b\},\{c\},\{d\},\{d'\},\{f\},\{r\}} \bar g^{a,e}(\{b\},\{c\},\{d\},\{d'\}) D^{ae,z}_{jklm}(\{b\},\{c\},\{d\},\{d'\},\{f\},\{r\})I_{jklm}(\{f\},\{r\}).
\end{multline}
and hence:
\begin{equation}
    \bra{\psi_m}\tilde{\mathcal{T}}^*_{N -z,\sigma_G}(n)[\ket{\psi_j}\bra{\psi_k}]\ket{\psi_l}
    =(1+v_{kl})(1+v_{jm})\bra{\psi_m}\tilde{\mathcal{B}}_{z,\sigma_G}(n)[\ket{\psi_j}\bra{\psi_k}]\ket{\psi_l}.
\end{equation}
As $|v_{jr}| \leq \epsilon_{\mathcal{B}} = O(\frac{n\beta \mu}{Nt})$ by Lemma~\ref{lemma:1}:
\begin{equation}
    \|\tilde{\mathcal{T}}^*_{\sigma_G}(n)[\mathbb{I}] - \tilde{\mathcal{B}}(n)[\sigma_G]\|_* \leq (2\epsilon_{\mathcal{B}} + \epsilon_{\mathcal{B}}^2)\|\tilde{\mathcal{B}}(n)[\sigma_G]\|_*
    \label{eq:Boutput}
\end{equation}
It remains to compare $\tilde{\mathcal{T}}^*_{\sigma_G}(n)[\mathbb{I}]$ and $\tilde{\mathcal{B}}^*_{\sigma_G}(n)[\mathbb{I}]$. Let $\tilde{\mathcal{T}}_z(n)$ denote $\tilde{\mathcal{B}}_z(n)$ with every $W_{abde}$ replaced by $W_{a-z,b-z,de}$, so that $\tilde{\mathcal{T}}^*_z(n)$ is its adjoint, and let $\tilde{\mathcal{T}}(n)=\sum_z\tilde{\mathcal{T}}_z(n)$. Exactly as in \eqref{eq:AAtildeclose1},
\begin{equation}
    \|\tilde{\mathcal{T}}^*_{\sigma_G}(n)[\mathbb{I}]-\tilde{\mathcal{B}}^*_{\sigma_G}(n)[\mathbb{I}]\|_*\leq Tr[\sigma_G]\,\|\tilde{\mathcal{T}}^{\dagger}(n)[\mathbb{I}]-\tilde{\mathcal{B}}^{\dagger}(n)[\mathbb{I}]\|\leq\|\tilde{\mathcal{T}}(n)-\tilde{\mathcal{B}}(n)\|_{\diamond},
    \label{eq:TBduality}
\end{equation}
so it suffices to compare the forward maps $\tilde{\mathcal{T}}(n)$ and $\tilde{\mathcal{B}}(n)$ in diamond norm. In order to use Lemma~\ref{lemma:3}, we must bound the difference in operator norm between $W_{abde}$ and $W_{a-z,b-z,de}$ for all $a,b,d,e$ and $|z|_{N} \leq (4n+6)\mu$. Therefore, we must appropriately deal with the modular cut in $\kappa_{abde}$. Due to the restricted support of $\tilde{\mathcal{T}}^*(n)$ and $\tilde{\mathcal{B}}^*(n)$ they only contain $W_{abde}$ and $W_{a-z,b-z,de}$ with $a,b,d,e, z$ such that $|\kappa_{abde}|,|\kappa_{a-z,b-z,de}| < \frac{N}{4}+(4n+6)\mu$. Therefore, we are free to modify $f$ outside of this region with no physical consequence, so we can replace $f$ by the continuous $N$-periodic extension $F$:
\begin{equation}
    F(\kappa) = \begin{cases}1& \text{ for }  -\frac{N}{2} \leq \kappa <  - \frac{3N}{8}\\
        \min(1,e^{-\beta \frac{2\pi}{Nt} \kappa}) & \text{ for } |\kappa|\leq \frac{3N}{8}\\
        e^{-\beta\frac{3\pi}{4t}} + \frac{1- e^{-\beta\frac{3\pi}{4t}}}{\frac{N}{8}}(\kappa -  \frac{3N}{8}) & \text{ for }  \frac{3N}{8} < \kappa <  \frac{N}{2}
    \end{cases}
\end{equation}
assuming $(4n+6)\mu < \frac{N}{8}$. Let $W^F_{abde}$ denote $W_{abde}$ with $f_{abde}$ replaced by $F(\kappa_{abde})$ which leaves $\tilde{\mathcal{B}}_z(n)$ and $\tilde{\mathcal{T}}_z(n)$ exactly unchanged. We use $W^F$ throughout the comparison below, and need the following lemma proved in Appendix~\ref{section:appendixTrunc}:
\begin{lemma}[$W$-drift]
\label{lemma:wdrift}
For all $a,b,d,e$ and all $|z|_{N} \leq (4n+6)\mu$,
\begin{equation}
    \|W^F_{a-z,b-z,de} - W^F_{abde}\| \leq 2\sqrt{\beta\frac{6\pi}{Nt}}\sqrt{(4n+6)\mu}.
\end{equation}
\end{lemma}
We define $\tilde Q_{s,z}$ as $\tilde Q_s$ with $W^F_{abde}$ substituted by $W^F_{a-z,b-z,de}$, and use $W^F$ in $\tilde Q_s$, $Q_s$ and $Q_{s,z}$. Therefore, by Lemma~\ref{lemma:3} we have:
\begin{equation}
    \|\tilde{\mathcal{Q}}_{S,z} - \tilde{\mathcal{Q}}_S\|_{\diamond}\leq \|\tilde{\mathcal{Q}}_{S,z} - \mathcal{Q}_{S,z}\|_{\diamond} +  \|\mathcal{Q}_{S,z} - \mathcal{Q}_S\|_{\diamond}+ \|\mathcal{Q}_{S} - \tilde{\mathcal{Q}}_S\|_{\diamond}\leq 12\sqrt{\beta\frac{6\pi}{Nt}}\sqrt{(4n+6)\mu}+24\epsilon_{\Delta}
\end{equation}
when $2\sqrt{\beta\frac{6\pi}{Nt}}\sqrt{(4n+6)\mu} \leq \frac{1}{4}$. Each stage of $\tilde{\mathcal{T}}_z(n)$ and $\tilde{\mathcal{B}}_z(n)$ has Kraus operators $\tilde U^\dagger D_s\tilde V$ with $\|\tilde U\|,\|\tilde V\|\leq1+2\epsilon_\Delta$ by Lemma~\ref{lemma:0} and each of the $2n+3$ stages has diamond norm at most $c=(1+2\epsilon_\Delta)^4$ and only the $n+1$ $Q$ stages differ. Hence, by the general form of Lemma~\ref{lemma:composition} with $k=2n+3$ and $c^{2n+2}\leq e^{2/3}<e$ under $12(2n+3)\epsilon_\Delta\leq1$, we can see that:
\begin{equation}
    \|\tilde{\mathcal{T}}_z(n) - \tilde{\mathcal{B}}_z(n)\|_{\diamond} \leq e(n+1)\left(12\sqrt{\beta\frac{6\pi}{Nt}}\sqrt{(4n+6)\mu}+24\epsilon_{\Delta}\right)
\end{equation}
Only $z$ with $|z|_{N} \leq (4n+6)\mu$ contribute, and there are at most $2(4n+6)\mu+1\leq 3(4n+6)\mu$ of these (as $\mu\geq1$), so
\begin{equation}
    \|\tilde{\mathcal{T}}(n) - \tilde{\mathcal{B}}(n)\|_{\diamond} \leq 36e(n+1)\sqrt{\beta\frac{6\pi}{Nt}}((4n+6)\mu)^{\frac{3}{2}}+72e(n+1)(4n+6)\mu\epsilon_{\Delta}
    \label{eq:Wdrifterror}
\end{equation}
Therefore, combining this with \eqref{eq:TBduality}, \eqref{eq:Boutput} and \eqref{eq:Bsrelation} via the triangle inequality, and using $\|\tilde{\mathcal{B}}(n)[\sigma_G]\|_*\leq\|\mathcal{B}(n)[\sigma_G]\|_*+24e(2n+3)\epsilon_{\Delta}$ from \eqref{eq:Bsrelation}, $\mathcal{B}(n)$ satisfies the hypothesis of Lemma~\ref{theorem:quantumapproximatedeatildedbalance} as:
\begin{equation}
    \|\mathcal{B}^*_{\sigma_G}(n)[\mathbb{I}] - \mathcal{B}(n)[\sigma_G]\|_* = O\left(\frac{(n+2)\beta\mu}{Nt} \right)\|\mathcal{B}(n)[\sigma_G]\|_* +O\left(\max\left((n+2)^2\mu\epsilon_{\Delta},(n+2)^{\frac{5}{2}}\sqrt{\frac{\beta }{Nt}}\mu^{\frac{3}{2}}\right)\right) 
\end{equation}

\section{Proofs of Lemmas}
\label{section:appendixTrunc}
\begin{lemma*}[\ref{lemma:0}]
% \label{lemma:0}
    Define $\widetilde{\text{BSI-QPE}}$ as given by:
    \begin{equation}
       \widetilde{\text{BSI-QPE}}=\sum_j \sum_{p,q} \tilde G^{pq}_j \otimes \ket{p}\bra{q} \otimes \ket{\psi_j}\bra{\psi_j}
    \end{equation}
    where 
    \begin{equation}
        \tilde G^{pq}_j = \begin{pmatrix}
        (-1)^{q-p}\tilde\alpha(q-p+\frac{N}{2\pi}E_jt) & \tilde\beta(q-p+\frac{N}{2\pi}E_jt)\\
        -\tilde\beta (q-p-\frac{N}{2\pi}E_jt) & (-1)^{q-p}\tilde\alpha(q-p-\frac{N}{2\pi}E_jt)
    \end{pmatrix}
    \end{equation}
    with $\tilde\alpha(x)=\hat\alpha(x)$, $\tilde\beta(x)=\hat\beta(x)$ for $|x|_{N}\leq\mu$ and zero otherwise, then,  
    \begin{equation}
        \|\widetilde{\text{BSI-QPE}}-\text{BSI-QPE}\| \leq 2\epsilon_{\Delta}= O\left(\frac{1}{\nu}e^{-\nu \mu}\right)
    \end{equation}
    where $\epsilon_{\Delta}$ is defined in \eqref{eq:decomposealpha3}.
\end{lemma*}
\begin{proof}
    As both operators are block diagonal in the Hamiltonian eigenbasis, we will have the operator norm given by the maximum over $j$:
    \begin{equation}
        \|\widetilde{\text{BSI-QPE}}-\text{BSI-QPE}\| = \max_j 
        \left\|\sum_{p,q}(\tilde G^{pq}_j - G^{pq}_j) \otimes \ket{p}\bra{q}\right\|
    \end{equation}
    As $\|U-V\| \leq \sqrt{\|U-V\|_{1\to1}\|U-V\|_{\infty\to\infty}}$ we can bound the spectral norm by bounding the $\|\cdot\|_{1\to1}$ (maximum absolute column sum) and $\|\cdot\|_{\infty\to\infty}$ (maximum absolute row sum) operator norms. First the $\|\cdot\|_{1\to1}$  operator norm:
    \begin{multline}
         \left\|\sum_{p,q}(\tilde G^{pq}_j - G^{pq}_j) \otimes \ket{p}\bra{q}\right\|_{1\to1} = \max_q \max\bigg(\sum_{p}|\tilde\alpha(q-p+\frac{N}{2\pi}E_jt)-\hat\alpha(q-p+\frac{N}{2\pi}E_jt)|+\\|\tilde\beta (q-p-\frac{N}{2\pi}E_jt)-\hat\beta (q-p-\frac{N}{2\pi}E_jt)|, \sum_{p}|\tilde\beta(q-p+\frac{N}{2\pi}E_jt)-\hat\beta(q-p+\frac{N}{2\pi}E_jt)|+\\|\tilde\alpha (q-p-\frac{N}{2\pi}E_jt)-\hat\alpha (q-p-\frac{N}{2\pi}E_jt)|\bigg)
    \end{multline}
    which simplifies to:
    \begin{multline}
         \left\|\sum_{p,q}(\tilde G^{pq}_j - G^{pq}_j) \otimes \ket{p}\bra{q}\right\|_{1\to1} = \max_q\max\bigg(\sum_{\substack{p\\|q-p+\frac{N}{2\pi}E_jt|_{N}>\mu}}|\hat\alpha(q-p+\frac{N}{2\pi}E_jt)|+\sum_{\substack{p\\|q-p-\frac{N}{2\pi}E_jt|_{N}>\mu}}|\hat\beta (q-p-\frac{N}{2\pi}E_jt)|,\\
         \sum_{\substack{p\\|q-p+\frac{N}{2\pi}E_jt|_{N}>\mu}}|\hat\beta(q-p+\frac{N}{2\pi}E_jt)|+\sum_{\substack{p\\|q-p-\frac{N}{2\pi}E_jt|_{N}>\mu}}|\hat\alpha (q-p-\frac{N}{2\pi}E_jt)|\bigg)
    \end{multline}
    Therefore, using \eqref{eq:scaling} and \eqref{eq:decomposealpha3}  we can see:
    \begin{equation}
        \left\|\sum_{p,q}(\tilde G^{pq}_j - G^{pq}_j) \otimes \ket{p}\bra{q}\right\|_{1\to1} \leq 2\epsilon_{\Delta} = O(\frac{1}{\nu}e^{-\nu \mu})
    \end{equation}
    A similar process gives us:
    \begin{equation}
        \left\|\sum_{p,q}(\tilde G^{pq}_j - G^{pq}_j) \otimes \ket{p}\bra{q}\right\|_{\infty\to\infty} \leq 2\epsilon_{\Delta} =O(\frac{1}{\nu}e^{-\nu \mu})
    \end{equation}
    Thus,
    \begin{equation}
        \|\widetilde{\text{BSI-QPE}}-\text{BSI-QPE}\| \leq 2\epsilon_{\Delta} = O\left(\frac{1}{\nu}e^{-\nu \mu}\right)
    \end{equation}
\end{proof}
\begin{lemma*}[\ref{lemma:3}]
% \label{lemma:3}
Let $\{D_s\}_{s\in S}$ be mutually orthogonal projectors satisfying $\sum_{s\in S} D_s \leq \mathbb{I}$. Given unitaries $U$ and $V$, define the measurement $K_s = U^{\dagger}D_s V$ and CP map $\mathcal{K}[\rho] = \sum_{s \in S} K_s\rho K_s^{\dagger}$. Further, given operators $\tilde U, \tilde V$, define approximate measurement operators $\tilde K_s = \tilde U^{\dagger}D_s \tilde V$ and the approximate map $\tilde{\mathcal{K}}[\rho] = \sum_{s \in S} \tilde K_s\rho \tilde K_s^{\dagger}$. If $\|U-\tilde U\| \leq \epsilon<\frac14$ and $\|V-\tilde V\| \leq \epsilon<\frac14$, then 
\begin{equation}
    \|\mathcal{K} - \tilde{\mathcal{K}}\|_{\diamond} \leq 6\epsilon
\end{equation}
\end{lemma*}
\begin{proof}
First, consider:
\begin{align}
    \mathcal{K}[\rho] - \tilde{\mathcal{K}}[\rho] &= \sum_{s \in S} K_s\rho K_s^{\dagger} - \tilde K_s\rho \tilde K_s^{\dagger} \\
    &= \sum_{s \in S} (K_s-\tilde K_s)\rho K_s^{\dagger} +\tilde K_s\rho (K_s^{\dagger}- \tilde K_s^{\dagger}) \label{eq:ksUdiff}
\end{align}
For $\Phi[\rho]=\sum_sR_{1,s}\rho R_{2,s}^\dagger$, define $V_i=\sum_s\ket{s}_S\otimes R_{i,s}$ for $i \in \{1,2\}$. Therefore, we can write $\Phi[\rho] = \operatorname{Tr}_S[V_1\rho V_2^{\dagger}]$. For any ancillary system $A$, the Schatten trace-norm duality gives:
\begin{align}
    \|(\Phi\otimes\operatorname{id}_A)[\rho]\|_* &= \sup_{\|X\|\leq 1} | \operatorname{Tr}[X (\operatorname{Tr}_S[(V_1\otimes \mathbb{I}_A) \rho (V_2^{\dagger}\otimes \mathbb{I}_A)]|\\
    &=\sup_{\|X\|\leq 1} | \operatorname{Tr}[(\mathbb{I}_S \otimes X ) (V_1\otimes \mathbb{I}_A) \rho (V_2^{\dagger}\otimes \mathbb{I}_A)]|\\
    &\leq\sup_{\|X\|\leq 1} \| (V_2^{\dagger}\otimes \mathbb{I}_A)(\mathbb{I}_S \otimes X ) (V_1\otimes \mathbb{I}_A) \rho \|_*\\
    &\leq \sup_{\|X\|\leq 1} \| V_2^{\dagger}\|\|X\| \|V_1\| \| \rho \|_* = \|V_1\|\| V_2\| \| \rho \|_*
\end{align}
where the last inequality comes from the generalised H\"older inequality of schatten norms. As $\|\Phi\|_{\diamond} = \sup_A\sup_{\|\rho\|_*\leq 1}\|(\Phi\otimes\operatorname{id}_A)[\rho]\|_*$ and
\(V_i^\dagger V_i=\sum_sR_{i,s}^\dagger R_{i,s}\), we have
\begin{equation}
\|\Phi\|_\diamond
\leq
\left\|\sum_{s\in S}R_{1,s}^\dagger R_{1,s}\right\|^{1/2}
\left\|\sum_{s\in S}R_{2,s}^\dagger R_{2,s}\right\|^{1/2}.
\end{equation}
Applying the above bound to Eq.~\eqref{eq:ksUdiff}
gives
\begin{equation}
\|\mathcal{K}-\tilde{\mathcal{K}}\|_\diamond
\leq
\left\|\sum_{s\in S}(K_s^\dagger - \tilde K_s^\dagger)(K_s - \tilde K_s)\right\|^{1/2}
\left(
\left\|\sum_{s\in S}K_s^\dagger K_s\right\|^{1/2}
+
\left\|\sum_{s\in S}\tilde K_s^\dagger\tilde K_s\right\|^{1/2}
\right).
\label{eq:row-column-bound}
\end{equation}
Let us evaluate the terms in this expression. First, we have
\begin{equation}
     \left\|\sum_{s\in S} K_s^{\dagger}  K_s \right\| = \left\|\sum_{s\in S}  V^{\dagger}D_sV \right\| \leq  \left\|V^{\dagger} V\right\| = 1
\end{equation}

Second, we have
\begin{align}
    \left\|\sum_{s\in S}\tilde K_s^{\dagger} \tilde K_s \right\|^{\frac{1}{2}} &= \left\|\sum_{s\in S} \tilde V^{\dagger} D_s \tilde U\tilde U^{\dagger}  D_s \tilde V\right\|^{\frac{1}{2}} \\
     &\leq \|\tilde U\|\left\|\sum_{s\in S} \tilde V^{\dagger} D_s  \tilde V\right\|^{\frac{1}{2}}\\
     &\leq \|\tilde U\| \|\tilde V^{\dagger}  \tilde V\|^{\frac{1}{2}}= \|\tilde U\|\|\tilde V\| \leq (1+\epsilon)^2\\
\end{align}
Finally let us consider:
\begin{align}
    K_s-\tilde K_s &= U^{\dagger}  D_s V-\tilde U^{\dagger}  D_s \tilde V\\
    &=  (U-\tilde U)^{\dagger}  D_s V+\tilde U^{\dagger}  D_s (V-\tilde V)\\
\end{align}
Let us define $E_U = U - \tilde U$, then:
\begin{align}
    \left\|\sum_{s\in S}(K_s^\dagger - \tilde K_s^\dagger)(K_s - \tilde K_s)\right\|^{\frac{1}{2}}&= \left\| \sum_{s \in S}\ket{s} \otimes ( E_U^{\dagger}  D_s V+\tilde U^{\dagger}  D_sE_V)\right\|\\
    &\leq  \left\| \sum_{s \in S} \ket{s} \otimes E_U^{\dagger}  D_s V \right\| + \left\| \sum_{s \in S} \ket{s} \otimes \tilde U^{\dagger}  D_sE_V\right\|\\
    &\leq \left\| \sum_{s \in S}  V^{\dagger}  D_sE_UE_U^{\dagger} D_s  V\right\|^{\frac{1}{2}} +  \left\| \sum_{s \in S} E_V^{\dagger}  D_s \tilde U \tilde U^{\dagger} D_s E_V \right\|^{\frac{1}{2}}\\
    &\leq \|E_U\| + \|E_V\|\|\tilde U\| \leq \epsilon(2+\epsilon)
\end{align}

Therefore,
\begin{equation}
    \|\mathcal{K} - \tilde{\mathcal{K}}\|_{\diamond} \leq \epsilon(2+\epsilon)(1+(1+\epsilon)^2)
\end{equation}
Using $\epsilon<\frac14$, we obtain $\|\mathcal{K} - \tilde {\mathcal{K}}\|_{\diamond} \leq 6\epsilon$.
\end{proof}
\begin{lemma*}[\ref{lemma:composition}]
Let $\mathcal{M}_1, \ldots, \mathcal{M}_k$ be completely positive trace-non-increasing maps and let $\tilde{\mathcal{M}}_1, \ldots, \tilde{\mathcal{M}}_k$ be completely positive maps with $\|\mathcal{M}_i - \tilde{\mathcal{M}}_i\|_{\diamond} \leq \varepsilon$ for all $i$. Then
\begin{equation}
    \|\mathcal{M}_k \circ \cdots \circ \mathcal{M}_1 - \tilde{\mathcal{M}}_k \circ \cdots \circ \tilde{\mathcal{M}}_1\|_{\diamond} \leq (1+\varepsilon)^k - 1 \leq e k \varepsilon,
\end{equation}
the last inequality holding whenever $k\varepsilon \leq 1$. More generally, if instead $\|\mathcal{M}_i\|_{\diamond}, \|\tilde{\mathcal{M}}_i\|_{\diamond} \leq c$ for some $c \geq 1$ and $\|\mathcal{M}_i - \tilde{\mathcal{M}}_i\|_{\diamond} \leq \varepsilon_i$, without assuming the $\mathcal{M}_i$ are trace-non-increasing, then
\begin{equation}
    \|\mathcal{M}_k \circ \cdots \circ \mathcal{M}_1 - \tilde{\mathcal{M}}_k \circ \cdots \circ \tilde{\mathcal{M}}_1\|_{\diamond} \leq c^{k-1}\sum_{i=1}^{k}\varepsilon_i .
\end{equation}
\end{lemma*}
\begin{proof}
    Writing the difference as a telescoping sum,
\begin{equation}
    \mathcal{M}_k \circ \cdots \circ \mathcal{M}_1 - \tilde{\mathcal{M}}_k \circ \cdots \circ \tilde{\mathcal{M}}_1 = \sum_{i=1}^{k} \mathcal{M}_k \circ \cdots \circ \mathcal{M}_{i+1} \circ (\mathcal{M}_i - \tilde{\mathcal{M}}_i) \circ \tilde{\mathcal{M}}_{i-1} \circ \cdots \circ \tilde{\mathcal{M}}_1,
\end{equation}
and using submultiplicativity of the diamond norm together with $\|\mathcal{M}_i\|_{\diamond} \leq 1$ (completely positive trace-non-increasing) and $\|\tilde{\mathcal{M}}_i\|_{\diamond} \leq \|\mathcal{M}_i\|_{\diamond} + \varepsilon \leq 1 + \varepsilon$, we obtain
\begin{equation}
    \|\mathcal{M}_k \circ \cdots \circ \mathcal{M}_1 - \tilde{\mathcal{M}}_k \circ \cdots \circ \tilde{\mathcal{M}}_1\|_{\diamond} \leq \sum_{i=1}^{k}\varepsilon(1+\varepsilon)^{i-1} = (1+\varepsilon)^k - 1.
\end{equation}
For $k\varepsilon \leq 1$ we have $(1+\varepsilon)^k \leq e^{k\varepsilon} \leq 1 + e k\varepsilon$, giving the final inequality. For the general statement, the $i$-th term of the same telescoping sum has diamond norm at most $c^{k-i}\,\varepsilon_i\,c^{i-1} = c^{k-1}\varepsilon_i$.
\end{proof}
\begin{lemma*}[\ref{lemma:summed-history-reversal}]
For $n\geq1$, let $\{x_i\}_{i=1}^{n+1}$ and $\{y_i\}_{i=1}^{n+1}$ denote
the compound labels
\begin{equation}
    y_i=(b_{2n+3-2i},d_{n+2-i}),\qquad x_i=(c_{2n+3-2i},d'_{n+2-i}),
\end{equation}
with $a,e$ fixed.  Choose $\theta_x\in[0,\pi/2]$ so that
\begin{equation}
    W_x=W_{bd}=\begin{pmatrix}
                \sqrt{1 - f_{abde}} & \sqrt{f_{abde}}\\
       \sqrt{f_{abde}} & -\sqrt{1 - f_{abde}}
    \end{pmatrix}= \begin{pmatrix}\cos\theta_x&\sin\theta_x\\
    \sin\theta_x&-\cos\theta_x\end{pmatrix},
\end{equation}
and set $\Delta_i=\theta_{y_i}-\theta_{x_i}$.  Then
\begin{equation}
    \bar g^{ae}(\{b\},\{c\},\{d\},\{d'\})
    =\frac12\cos\Delta_1\cos\Delta_{n+1}
      \prod_{i=1}^{n}\cos(\Delta_i-\Delta_{i+1}).
\end{equation}
Consequently,
\begin{equation}
    \bar g^{ae}(\{b\},\{c\},\{d\},\{d'\})
    =\bar g^{ae}(\{c\}^R,\{b\}^R,\{d'\}^R,\{d\}^R).
    \label{eq:PRR}
\end{equation}
\end{lemma*}
\begin{proof}
Put $\ket{v_s(x)}=W_x\ket{s}$, so
$R_s^x=\ket{v_s(x)}\bra{v_s(x)}$. Therefore, we can rewrite $\bar g^{ae}$ for $n\geq 1$ \eqref{eq:gbar-definition} as:
\begin{multline}
    \frac{1}{2}\sum_{s_0,\ldots,s_n} \braket{0|v_{s_0}(x_1)}\braket{v_{s_0}(x_1)|v_{s_1}(x_2)}\ldots \braket{v_{s_{n-1}}(x_n)|v_{s_n}(x_{n+1})}\braket{v_{s_n}(x_{n+1})|v_{s_n}(y_{n+1})}\\\braket{v_{s_n}(y_{n+1})|v_{s_{n-1}}(y_n)}\ldots\braket{v_{s_1}(y_2)|v_{s_0}(y_1)}\braket{v_{s_0}(y_1)|0}
\end{multline}
We can reorder this to get:
\begin{multline}
    \frac{1}{2}\sum_{s_0,\ldots,s_n} \braket{0|v_{s_0}(x_1)}\braket{v_{s_0}(y_1)|0}\braket{v_{s_0}(x_1)|v_{s_1}(x_2)}\braket{v_{s_1}(y_2)|v_{s_0}(y_1)}\ldots\\ \braket{v_{s_{n-1}}(x_{n})|v_{s_{n}}(x_{n+1})}\braket{v_{s_{n}}(y_{n+1})|v_{s_{n-1}}(y_{n})}\braket{v_{s_{n}}(x_{n+1})|v_{s_{n}}(y_{n+1})}
\end{multline}
For consecutive labels define
the two-state transfer matrix
\begin{equation}
    (T_i)_{uv}=\braket{v_u(x_i)|v_v(x_{i+1})}
    \braket{v_v(y_{i+1})|v_u(y_i)}.
\end{equation}
Therefore, we can rewrite $\bar g$ as:
\begin{equation}
   \frac{1}{2}\sum_{s_0,\ldots,s_n} h_{s_0}(T_1)_{s_0s_1}(T_2)_{s_1s_2}\ldots (T_n)_{s_{n-1}s_{n}} c_{s_n}
\end{equation}
where $h_s = \braket{0|v_{s}(x_1)}\braket{v_{s}(y_1)|0}$ and $c_s =\braket{v_{s}(x_{n+1})|v_{s}(y_{n+1})}$. Hence, we can write $\bar g$ as simple matrix multiplication:
\begin{equation}
    \bar g(\{b\},\{c\},\{d\},\{d'\}) = \frac{1}{2}hT_1T_2T_3\ldots T_nc
\end{equation}
As $W_xW_y = \begin{pmatrix}
    \cos (\theta_x - \theta_y) & - \sin (\theta_x - \theta_y)\\
     \sin (\theta_x - \theta_y) & \cos (\theta_x - \theta_y)
\end{pmatrix}$ we have:
\begin{align}
    (T_i)_{u v} &=  \begin{pmatrix}
    \cos (\theta_{x_i} - \theta_{x_{i+1}}) & - \sin (\theta_{x_i} - \theta_{x_{i+1}})\\
     \sin (\theta_{x_i} - \theta_{x_{i+1}}) & \cos (\theta_{x_i} - \theta_{x_{i+1}})
\end{pmatrix}_{uv} \begin{pmatrix}
    \cos (\theta_{y_{i+1}} - \theta_{y_{i}}) & - \sin (\theta_{y_{i+1}} - \theta_{y_{i}})\\
     \sin (\theta_{y_{i+1}} - \theta_{y_{i}}) & \cos (\theta_{y_{i+1}} - \theta_{y_{i}}) 
\end{pmatrix}_{vu}\\
&=\begin{pmatrix}
    \cos (\theta_{x_i} - \theta_{x_{i+1}})\cos (\theta_{y_{i+1}} - \theta_{y_{i}}) & - \sin (\theta_{x_i} - \theta_{x_{i+1}})\sin (\theta_{y_{i+1}} - \theta_{y_{i}})\\
     -\sin (\theta_{x_i} - \theta_{x_{i+1}})\sin (\theta_{y_{i+1}} - \theta_{y_{i}}) & \cos (\theta_{x_i} - \theta_{x_{i+1}}) \cos (\theta_{y_{i+1}} - \theta_{y_{i}}) 
\end{pmatrix}_{uv}\\
&=  \cos (\theta_{y_i} - \theta_{y_{i+1}})\cos (\theta_{x_{i}} - \theta_{x_{i+1}}) \mathbb{I}_{uv} + \sin (\theta_{y_i} - \theta_{y_{i+1}})\sin (\theta_{x_{i}} - \theta_{x_{i+1}}) X_{uv}
\end{align}
We can also clearly see that $c = \cos (\theta_{x_{n+1}} -\theta_{y_{n+1}}) \begin{pmatrix}
    1\\1
\end{pmatrix} =  \cos(\Delta_{n+1})\begin{pmatrix}
    1\\1
\end{pmatrix} $, and we have that $\begin{pmatrix}
    1\\1
\end{pmatrix}$ is an eigenvector of all $T_i$ as $T_i = \alpha_i \mathbb{I} + \beta_i X$. The corresponding eigenvalue for each $T_i$ is given by:
\begin{align}
    \alpha_i+\beta_i&=\cos (\theta_{y_i} - \theta_{y_{i+1}})\cos (\theta_{x_{i}} - \theta_{x_{i+1}}) + \sin (\theta_{y_i} - \theta_{y_{i+1}})\sin (\theta_{x_{i}} - \theta_{x_{i+1}}) \\
    &= \cos(\theta_{y_i} - \theta_{y_{i+1}} - (\theta_{x_{i}} - \theta_{x_{i+1}}))\\
    &= \cos\!\left(\theta_{y_i} - \theta_{x_i} - (\theta_{y_{i+1}} - \theta_{x_{i+1}})\right)\\
    &= \cos(\Delta_i - \Delta_{i+1})
\end{align}
Therefore, $T_1T_2T_3\ldots T_nc = \cos(\Delta_{n+1})\prod_{i=1}^{n}\cos(\Delta_i - \Delta_{i+1})\begin{pmatrix}
    1\\1
\end{pmatrix}$. As $h = \begin{pmatrix}
    \cos(\theta_{x_1})\cos(\theta_{y_1})& \sin(\theta_{x_1})\sin(\theta_{y_1})
\end{pmatrix}$, we have $h\begin{pmatrix}
    1\\1
\end{pmatrix}= \cos(\theta_{x_1})\cos(\theta_{y_1}) + \sin(\theta_{x_1})\sin(\theta_{y_1}) = \cos(\theta_{x_1}-\theta_{y_1}) = \cos(\Delta_1)$. Therefore,
\begin{equation}
    \bar g(\{b\},\{c\},\{d\},\{d'\}) =\frac{1}{2}\cos(\Delta_1)\cos(\Delta_{n+1})\prod_{i=1}^{n}\cos(\Delta_i - \Delta_{i+1})
\end{equation}
  Exchanging $b,d'$ ($x$) with $c,d$ ($y$) maps $\Delta_i\mapsto-\Delta_i$, and $\bar g$ is preserved because cosine is even. Reversing $\{b\}$, $\{c\}$, $\{d\}$ and $\{d'\}$ maps $\Delta_i\mapsto\Delta_{n+2-i}$. As $\frac{1}{2}\cos(\Delta_{n+1})\cos(\Delta_{1})\prod_{i=1}^{n}\cos(\Delta_{n+2-i} - \Delta_{n+1-i}) = \frac{1}{2}\cos(\Delta_1)\cos(\Delta_{n+1})\prod_{i=1}^{n}\cos(\Delta_i - \Delta_{i+1})$, the quantity $\bar g$ is preserved and Eq.~\eqref{eq:PRR} holds.
\end{proof}
\begin{lemma*}[\ref{lemma:wdrift} $W$-drift]
% \label{lemma:wdrift}
For all $a,b,d,e$ and all $|z|_{N} \leq (4n+6)\mu$,
\begin{equation}
    \|W^F_{a-z,b-z,de} - W^F_{abde}\| \leq 2\sqrt{\beta\frac{6\pi}{Nt}}\sqrt{(4n+6)\mu}.
\end{equation}
\end{lemma*}
\begin{proof}
    By definition we have
\begin{equation}
     W^F_{a-z,b-z,de} = \begin{pmatrix}
        \sqrt{1 - F(\kappa_{a-z,b-z,de})} & \sqrt{F(\kappa_{a-z,b-z,de})}\\
       \sqrt{F(\kappa_{a-z,b-z,de})} & -\sqrt{1 - F(\kappa_{a-z,b-z,de})}
    \end{pmatrix}
\end{equation}
where:
\begin{equation}
    F(\kappa) = \begin{cases}1& \text{ for }  -\frac{N}{2} \leq \kappa <  -\frac{3N}{8}\\
        \min(1,e^{-\beta \frac{2\pi}{Nt} \kappa}) & \text{ for } |\kappa|\leq \frac{3N}{8}\\
        e^{-\beta\frac{3\pi}{4t}} + \frac{1- e^{-\beta\frac{3\pi}{4t}}}{\frac{N}{8}}(\kappa - \frac{3N}{8}) & \text{ for }  \frac{3N}{8} < \kappa <  \frac{N}{2}
    \end{cases}
\end{equation}
and
\begin{equation}
    \kappa_{a-z,b-z,de} \equiv (-1)^d(b-z) - (-1)^e(a-z) \mod N \text{ for } -\frac{N}{2} \leq \kappa_{a-z,b-z,de} < \frac{N}{2}
\end{equation}
As $|z|_{N} \leq (4n+6)\mu< \frac{N}{8}$  we have:
\begin{equation}
    |\kappa_{a-z,b-z,de} - \kappa_{abde}|_{N} \leq 2(4n+6)\mu
\end{equation}
The function $F$ is continuous and $N$-periodic, so its Lipschitz constant with respect to the circular distance $|\cdot|_{N}$ is given by its largest gradient. For $|\kappa|\leq \frac{3N}{8}$ the largest gradient is given by $\frac{2\pi\beta}{Nt}$, whereas for $ \frac{3N}{8} < \kappa <  \frac{N}{2}$ the gradient is given by $\frac{1- e^{-\beta\frac{3\pi}{4t}}}{\frac{N}{8}} \leq \frac{6\pi\beta}{Nt}$. Therefore,
\begin{equation}
     |F(\kappa)-F(\kappa')|\leq \frac{6\pi\beta}{Nt}|\kappa - \kappa'|_{N} \qquad \forall \quad \kappa, \kappa'
\end{equation}
Hence,
\begin{equation}
    |F(\kappa_{a-z,b-z,de})-F(\kappa_{abde})|\leq 12(4n+6)\frac{\pi\mu\beta}{Nt}
\end{equation}
We have
\begin{equation}
    W^F_{a-z,b-z,de}-W^F_{abde}=\begin{pmatrix} p & q\\ q & -p\end{pmatrix}
\end{equation}
where $p=\sqrt{1-F(\kappa_{a-z,b-z,de})}-\sqrt{1-F(\kappa_{abde})}$ and $q=\sqrt{F(\kappa_{a-z,b-z,de})}-\sqrt{F(\kappa_{abde})}$. This matrix is real symmetric and squares to $(p^2+q^2)\mathbb{I}$, so its operator norm is $\sqrt{p^2+q^2}$. Since $|\sqrt{x}-\sqrt{y}|\le\sqrt{|x-y|}$ for all $x,y\ge 0$, we have $|p|\leq \sqrt{12(4n+6)\frac{\pi\mu\beta}{Nt}}$ and $|q|\leq \sqrt{12(4n+6)\frac{\pi\mu\beta}{Nt}}$. Therefore
\begin{equation}
    \|W^F_{a-z,b-z,de}-W^F_{abde}\|\le\sqrt{24(4n+6)\frac{\pi\mu\beta}{Nt}}
    =2\sqrt{\beta\frac{6\pi}{Nt}}\sqrt{(4n+6)\mu}
\end{equation}
\end{proof}
